\section{A completely explicit rational Gaussian seed}
\label{gaussian:sec}

We now instantiate Section~\ref{amp:sec}.  The parameters are deliberately
conservative: their purpose is to give a reproducible numerical theorem
using rational inequalities, with no unspecified ``sufficiently large''
choice.  The resulting exponent is very small.  The construction complements
the small exact witness and the codimension classification rather than
providing a practical Boolean example.

\subsection{Intervals and selector}

Set
\begin{equation}
 \begin{aligned}
 s&=2^{13},\qquad m=s^2+1=2^{26}+1,\qquad R=9,\\
 c&=\frac13,\qquad h=\frac9{s^3}=\frac9{2^{39}}.
 \end{aligned}
 \label{gaussian:parameters}
\end{equation}
For $0\le i<m$, take
\begin{equation}
 a_i=c-9+\frac{18i}{m-1},\qquad I_i=[a_i-h,a_i+h].
 \label{gaussian:intervals}
\end{equation}
For $|t|\le9$, the selector chooses a center nearest to $t+c$, with
ties resolved by the smaller index.  For $|t|>9$, it chooses the middle
index, whose center is $c$.  Every endpoint and cut point is rational.
Each selector cell is an interval except for the middle cell, which is
a union of at most three intervals.  Thus $\ell=3$ is valid in
Theorem~\ref{amp:transfer}.

Let $q_i,\mu_i,s_i$ be the Gaussian interval probabilities, means, and
variances from \eqref{amp:interval-parameters}, and write
$Q=\prod_iq_i$, $S=\sum_i s_i$, and $J=J_G$.  Define
\[
 r_0=\frac9{s^2},\qquad e_0=r_0+h.
\]
The elementary interval estimates are
\begin{equation}
 |\mu_i-a_i|\le h,\qquad 0\le s_i\le h^2,\qquad
 S\le mh^2=\frac{81}{s^4}(1+s^{-2}).
 \label{gaussian:interval-estimates}
\end{equation}
On $|t|\le9$, nearest-center selection additionally gives
$|t+c-\mu_{\iota(t)}|\le e_0$.

\begin{lemma}[Certified Gaussian gain]
\label{gaussian:gain}
The parameters \eqref{gaussian:parameters}--\eqref{gaussian:intervals}
satisfy
\begin{equation}
 J_{\rm central}<\frac3{40},\qquad
 J_{\rm tail}<\frac1{1000},\qquad
 \Gamma:=c^2+S-J>\frac1{32}.
 \label{gaussian:gain-bounds}
\end{equation}
Moreover every interval is contained in $[-10,10]$, and
\begin{equation}
 q:=\min_iq_i>
 q_0:=\frac1{2^{38}3^{49}}>2^{-116}.
 \label{gaussian:q-bound}
\end{equation}
\end{lemma}

\begin{proof}
Write $\phi(t)=(2\pi)^{-1/2}e^{-t^2/2}$.  If $|t|\le9$ and
$z\in I_{\iota(t)}$, then $z-t=c+\xi$ with $|\xi|\le e_0$.
Consequently
\[
 \frac{z^2-t^2}{2}
 \le ct+\frac{c^2}{2}+e_0(9+c)+\frac{e_0^2}{2}
 <ct+\frac1{16}.
\]
The last inequality is a rational comparison after substituting
\eqref{gaussian:parameters}; it is also among the checks in the
constant certificate below.  Taking the minimum density on the selected
interval gives
\begin{equation}
 \frac{\phi(t)}{q_{\iota(t)}}
 \le\frac{e^{1/16}e^{ct}}{2h}
 \qquad(|t|\le9).
 \label{gaussian:density-ratio}
\end{equation}
Inserting \eqref{gaussian:interval-estimates} into
\eqref{amp:J} and integrating $e^{ct}$, we obtain
\begin{align}
 J_{\rm central}
 &\le\frac{3(e_0^2+mh^2)}{h}e^{1/16}\sinh3\notag\\
 &=\frac{54}{s}(1+s^{-1}+s^{-2})e^{1/16}\sinh3.
 \label{gaussian:central}
\end{align}
The elementary series bound $e^x<(1-x)^{-1}$ for $0<x<1$
gives $e^{1/16}<16/15$.  Also $e<11/4$ implies
$\sinh3<e^3/2<21/2$.  Thus the last expression is less than
\[
 \frac{54}{8192}(1+8192^{-1}+8192^{-2})\frac{16}{15}\frac{21}{2}
 <\frac3{40}.
\]
For an entirely rational justification of the bound on $e$, sum its
series through $1/3!$ and bound the remaining factorial tail by
$(1/24)\sum_{j\ge0}5^{-j}=5/96$.  This gives
$e<87/32<11/4$; the strict lower bound $e>5/2$ follows from the
first three terms and the positive remainder.

On $|t|>9$ the fallback interval has center $c$, and
$c+h<1/2$.  Its probability is at least $2h\phi(1/2)$.
Since $h<1$, $S<1$, and $|t|\ge9$,
\[
 (t+c-\mu_{\rm middle})^2+\sum_{j\ne\rm middle}s_j
 \le(|t|+h)^2+S\le2t^2.
\]
Integration by parts and the Gaussian tail bound give
\[
 \int_9^\infty t^2e^{-t^2/2}\,dt
 \le\left(9+\frac19\right)e^{-81/2}.
\]
Using both tails therefore yields
\begin{equation}
 J_{\rm tail}\le\frac{164}{9h}e^{-323/8}
 <\frac{55}{3h}e^{-323/8}
 <\frac{55\,2^{79}}{27\,5^{40}}
 <\frac1{1000}.
 \label{gaussian:tail}
\end{equation}
Here $323/8>40$ and $e>5/2$ justify the penultimate exponential
comparison.  In particular,
\[
 \Gamma\ge c^2-J
 >\frac19-\frac3{40}-\frac1{1000}
 =\frac{79}{2250}>\frac1{32}.
\]

Finally, $|a_i|+h<10$ for every $i$.  Since $\sqrt{2\pi}<3$,
\[
 q_i\ge2h\min_{z\in I_i}\phi(z)
 >\frac{2h}{3}e^{-50}
 >\frac{2h}{3}\,3^{-50}
 =\frac1{2^{38}3^{49}}.
\]
The last strict inequality uses $e<3$.  The integer inequality
$3^{49}<2^{78}$ now proves \eqref{gaussian:q-bound}.
\end{proof}

\subsection{A fixed averaging count and an explicit exponent}

\begin{theorem}[Fully numerical quantitative instance]
\label{gaussian:explicit}
For the rational reporting rule above, set
\begin{equation}
 M=2^{1123},\qquad
 \eta_0=2^{-(116m+7)},\qquad
 \varepsilon_0=2^{-8\,000\,000\,000}.
 \label{gaussian:numerical-choices}
\end{equation}
There is a sequence of exactly specified finite Boolean functions $f_k$
on $((m+1)M)^k$ bits, with
\begin{equation}
 \begin{aligned}
 \deg(f_k)&\le(mM)^k,\\
 L(f_k)&\ge6^{-1/2}\bigl(M(m+\eta_0)\bigr)^{k/2}
       \ge6^{-1/2}\deg(f_k)^{1/2+\varepsilon_0}.
 \end{aligned}
 \label{gaussian:explicit-growth}
\end{equation}
Their degrees tend to infinity.  For every degree budget $D\ge mM$,
the functions in Corollary~\ref{amp:dimension} can be chosen on at most
\begin{equation}
 D^{1+2^{-35}}
 \label{gaussian:budget-dimension}
\end{equation}
bits.  In terms of the actual degree $d=\deg(f_k)$ of an iterate,
\begin{equation}
 N_k\le(\sqrt6\,d)^{2+2^{-34}}.
 \label{gaussian:actual-dimension}
\end{equation}
\end{theorem}

\begin{proof}
We verify the transfer directly with
\[
 K=6,\quad B=10,\quad\ell=3,\quad C=|c|+B<11,
 \qquad\delta=2^{-560}.
\]
Using $m<2^{27}$, $q>2^{-116}$, $\sqrt6<3$, and
$6^{1/4}<2$, the constant in \eqref{amp:A} is bounded by the
explicit integer
\begin{align}
 A<&\ 3600m\,2^{116}
     +2m\,2^{116}(1152+600m)\notag\\
    &\quad +(6336+7600m)\,2^{232}\notag\\
   <&\ 2^{274}.
 \label{gaussian:A-bound}
\end{align}
Consequently $A\sqrt\delta<2^{-6}=1/64<\Gamma/2$.
Also $\delta<q/(4m)$, so the proof of
Theorem~\ref{amp:transfer} applies at this value of $\delta$.
Since
\[
 \sqrt{6/M}=\sqrt6\,2^{-1123/2}<2^{-560},
\]
the same count $M$ works at every stage of
Theorem~\ref{amp:fixed-count}.  Lemma~\ref{gaussian:gain} further gives
\[
 \frac{Q\Gamma}{4}>
 \frac{(2^{-116})^m}{128}=\eta_0.
\]
Thus $\eta_0$ is an admissible lower gain, even though the exact
Gaussian interval probabilities are not evaluated.

Write
\[
 \varepsilon=
 \frac{\log(1+\eta_0/m)}{2\log(mM)}.
\]
For $0<x<1$, $\log(1+x)>x/2$.  Moreover
$\log(mM)<1150$ and $4600m<2^{40}$.  Hence
\begin{equation}
 \begin{aligned}
 \varepsilon&>\frac{\eta_0}{4600m}>2^{-(116m+47)}\\
             &=2^{-7\,784\,628\,387}>\varepsilon_0.
 \end{aligned}
 \label{gaussian:epsilon}
\end{equation}
The growth assertion follows from Theorem~\ref{amp:fixed-count}.

For the dimension estimate, use $\log(1+1/m)<1/m$ and
$\log2>1/2$ to obtain
\[
 \theta-1
 =\frac{\log(1+1/m)}{\log(mM)}
 <\frac2{1123m}<2^{-35}.
\]
The last step is the integer inequality $1123m>2^{36}$.
This proves \eqref{gaussian:budget-dimension}.
Since $v_*=M(m+\eta_0)>mM=d_*$, the exponent in
\eqref{amp:actual-degree} satisfies
$\beta<2\theta<2+2^{-34}$, proving
\eqref{gaussian:actual-dimension}.
\end{proof}

\subsection{What the finite certificate verifies}

The accompanying script \texttt{code/verify\_constants.py} uses Python
integer and rational arithmetic.  It checks the density-exponent margin,
the rational bounds in \eqref{gaussian:central} and
\eqref{gaussian:tail}, the variance and endpoint bounds, the probability
comparison $3^{49}<2^{78}$, the integer upper bound in
\eqref{gaussian:A-bound}, the averaging inequality, and the exponent and
dimension comparisons.  Its recorded output is
\texttt{data/constant\_certificate.json}.  The analytic implications of
those rational comparisons are proved above; the script does not approximate
Gaussian probabilities or certify a numerical normal-distribution table.

The two very small dyadic constants in
\eqref{gaussian:numerical-choices} are stored by their integer exponents.
Neither the verification nor the construction of the certificate allocates
an integer with billions of bits.  The Boolean construction is nevertheless
finite and algorithmically specified.  Its size precludes using the displayed
parameters to generate a truth table in practice.
